\hypertarget{chapter-15-computation-inside-the-document}{%
\chapter{Computation Inside the Document}\label{chapter-15-computation-inside-the-document}}

The previous chapter treated figures as derived data rather than attachments. This chapter extends the same principle to computation more broadly: code listings that are verified rather than merely typed, simulation output that regenerates automatically, and a build process where every commit produces a document that is guaranteed to reflect the current state of the underlying computation, not a snapshot someone remembered to update the last time they happened to rerun it by hand.

\hypertarget{verified-code-listings-with-minted}{%
\section{\texorpdfstring{Verified Code Listings with \texttt{minted}}{Verified Code Listings with minted}}\label{verified-code-listings-with-minted}}

Most LaTeX documents that include source code do so with \texttt{verbatim} or \texttt{listings}, both of which typeset whatever text is given to them without checking that it is valid, runnable code in the language it claims to be. \texttt{minted}, built on the external Pygments syntax highlighter, provides considerably better highlighting across a wide range of languages, but its more important property for a technical book is what it makes possible in combination with a build pipeline: because \texttt{minted} can also execute code through Pygments' broader ecosystem tooling, or be paired with a build step that runs the listed code separately and compares its output against what the document claims that output to be, a code listing can be checked rather than merely displayed.

This distinction matters most for a book, like this one, whose worked examples are meant to actually work. A hand-typed code listing that once ran correctly, at the time it was written, can silently drift out of sync with the library or language version it was written against, and nothing about a plain \texttt{verbatim} block catches this drift; the listing continues to render exactly as typed regardless of whether it would still run. A listing pulled from an actual source file checked into the project, included with \texttt{\textbackslash{}inputminted\{python\}\{scripts/example.py\}} rather than retyped inline, cannot drift from what that file contains, and if the project's build process also runs that file as part of its test suite, a change that breaks the example is caught at build time rather than silently shipped in the next PDF.

\hypertarget{calling-out-to-external-languages-at-build-time}{%
\section{Calling Out to External Languages at Build Time}\label{calling-out-to-external-languages-at-build-time}}

For simulation results, numerical examples, or any computed value a book wants to report precisely, retyping a number by hand from wherever it was computed is the same category of error as hand-transcribing a table from a CSV file, covered in the previous chapter: the number in the document can silently diverge from what the underlying computation actually produces the moment either one is updated without the other. The more robust pattern, extending the Makefile-driven approach from Chapter 14, is a build step that runs the relevant Python, Rust, MATLAB, R, or Julia code, writes its output, whether a number, a table, or a figure, to a file the document then reads, and rebuilding the document means rerunning that computation first.

For a single computed value referenced inline in prose, this can be as simple as a script that writes a small \texttt{.tex} file defining a macro (\texttt{\textbackslash{}newcommand\{\textbackslash{}resultvalue\}\{4.7\}}), included into the main document, so that \texttt{\textbackslash{}resultvalue} in the prose always reflects the most recently computed result rather than a number someone typed once and never revisited. For more substantial output, a full table or figure, the same CSV-based pipeline from the previous chapter applies directly, with the computation itself, not merely its formatting, as the first step of the pipeline rather than something that happened once, off to the side, before the document's own build process began.

\hypertarget{the-document-as-the-final-stage-of-a-computation}{%
\section{The Document as the Final Stage of a Computation}\label{the-document-as-the-final-stage-of-a-computation}}

Adopting this pattern consistently changes what a ``build'' of the document actually means. Rather than a document being written up after a computation is finished, with numbers and figures copied over by hand at whatever point the writing catches up to the results, the document's own build process becomes the last stage of the computation itself: run the analysis, generate the figures and computed values, then typeset the document that presents them, all as one sequence triggered by a single command.

This has a consequence worth stating directly, because it changes how a document should be trusted. A document built this way cannot present a stale result without the build itself failing to actually rerun the computation, which is a detectable, fixable failure of the build process rather than a silent, undetectable failure of an author's memory to update a hand-copied number. A reviewer, or a reader building the project from source months after publication, gets the same guarantee: the numbers in front of them are what the checked-in code actually produces, right now, on this machine, not what it produced once, on the author's machine, at some earlier point that may no longer match the current state of the code.

\hypertarget{continuous-integration-a-fresh-pdf-on-every-commit}{%
\section{Continuous Integration: A Fresh PDF on Every Commit}\label{continuous-integration-a-fresh-pdf-on-every-commit}}

The natural extension of this discipline is running the entire pipeline, computation and typesetting together, automatically on every commit, rather than only when an author remembers to run it locally before sharing a draft. A continuous integration workflow, using GitHub Actions or a comparable service, checked into the repository alongside the document's own source, can be configured to install the project's dependencies (the TeX distribution, the plotting and simulation libraries, any language runtimes the build scripts require), run the full build, computation and all, and publish the resulting PDF as a build artifact attached to that commit.

This closes the loop the earlier sections in this chapter opened. A broken example, a computation that no longer runs against a current library version, or a figure-generating script that has drifted out of sync with the data it expects, is caught automatically, on every commit, rather than discovered by a reader months later or by the author themselves the next time they happen to rebuild the project from scratch. For a book-length project under active development for years, where memory of exactly how each figure was originally generated fades long before the project is finished, this is not a convenience; it is what keeps a large, computationally grounded document trustworthy over a timescale longer than any single person's memory of their own build process can be relied on to cover.

Part VI closes here, with graphics and computation both treated as part of the document's own build rather than as external inputs prepared separately and pasted in. Part VII turns to document engineering proper: custom document classes, professional title pages and front matter, publishing pipelines, and debugging a large project when, despite all of this discipline, something still goes wrong.
